\clearpage\SourceLeaf{77}\NoSourceHead
\begin{center}
{\Large ANHANG.}\par\bigskip
{\Large ÜBER EINEN SATZ VON GAUSS.}\par\medskip
Von M. A. \textsc{Stern.}\par\medskip
\rule{.65\textwidth}{.3pt}\par
{\small Nachrichten von der Königlichen Gesellschaft der Wissenschaften\par
und der G. A.-Universität zu Göttingen, 1869, S. 330—334.}\par
\rule{.65\textwidth}{.3pt}
\end{center}
\noindent\hspace*{\parindent}Unter den wissenschaftlichen Papieren meines verstorbenen Freundes Dr. M. \textsc{Reiss}, deren Durchsicht ich im Auftrage seiner Familie unternommen habe, findet sich ein aus 32 Quartseiten bestehendes Heft, welches ohne Zweifel gänzlich von \textsc{Gauss} geschrieben ist. Es besteht fast nur aus Formeln und Rechnungen, selten sind einzelne Worte eingestreut. Die Entstehung dieses Heftes ist leicht zu erklären. \textsc{Reiss} hat nemlich von Ostern 1823 bis Ostern 1825 hier in Göttingen studirt und bei \textsc{Gauss} Privatissima gehört. Das erwähnte Heft enthält die Entwicklungen, welche \textsc{Gauss} während des mündlichen Vortrages niederschrieb. Es ergibt sich dies deutlich aus einem andern von \textsc{Reiss} geschriebenen und unter seinen Papieren befindlichen Hefte, welches die Aufschrift hat: \textit{Ausarbeitung des in dem Vortrage des Herrn Hofrath Gauss Enthaltenen, angefangen den 8. Nov. 1824.} Dieses \textsc{Reiss}sche Heft nebst einigen dazu gehörenden losen Blättern enthält, wie die Übereinstimmung der Formeln zeigt, die Ausarbeitung dessen, was auf den ersten 10 Seiten des von \textsc{Gauss} geschriebenen Heftes vorkommt und z. B. den Beweis des \textsc{Harriot}schen Lehrsatzes, welchen \textsc{Gauss} später in dem \textsc{Crelle}schen Journal f. d. Mathematik [\textbf{3}, 1828, S. 1, Werke III, S. 65] bekannt gemacht hat. Mehr habe ich nicht aufgefunden. An einer Stelle des \textsc{Reiss}schen Heftes wird eine von \textsc{Gauss} im Sommer 1824 ersonnene Methode erwähnt. Man kann hieraus in Verbindung mit dem oben erwähnten Datum schliessen, dass das von \textsc{Gauss} geschriebene Heft im Laufe des Wintersemesters 1824—1825 entstanden ist.
\newpage\SourceLeaf{78}\ArithmeticHead{71}{R}{M. A. STERN, ÜBER EINEN SATZ VON GAUSS.}
Nur an einer Stelle dieses Heftes findet sich eine grössere Zahl zusammenhängender Worte und diese enthalten einen Satz aus der höheren Arithmetik. Die Worte lauten:
\begin{center}\begin{tabular}{@{}l@{}}
\(a\) ganze (pos. oder neg.) Zahl von der Form \(4k+1\)\\
\(n\) beliebige ganze positive Zahl ungerade\\
\(a\) und \(n\) sollen keinen gemeinschaftlichen Theiler haben\\
\(f\) alle ungeraden Zahlen \(1,3,5,7,\ldots,n-2\)\\
\(q\) ganzer Theil des Quotienten \(\frac{af}n\).
\end{tabular}\end{center}
\noindent\hspace*{\parindent}Unter allen \(q\) finden sich ebensoviele von der Form \(4k+2\) wie von der Form \(4k+3\).

Hierauf folgen Formeln, welche offenbar in keiner Beziehung zu diesen Worten stehen. Dann aber kommen Andeutungen eines Beweises des in diesen Worten enthaltenen Satzes, welche man leicht ergänzen kann. Ich lasse hier diesen Beweis folgen, indem ich bemerke, dass ich das von mir hinzugefügte in Klammern eingeschlossen habe; alles Übrige steht in dem Hefte.

[Es bezeichne]
\begin{center}\begin{tabular}{@{}l@{}}
\(f\) die Zahlen \(1,3,5,7,9,\ldots,n-2\)\\
\(g\) die Zahlen \(2,4,6,8,10,\ldots,n-1\)\\
\(h\) die Zahlen \(1,2,3,4,\ldots,\frac12(n-1)\)\\
\(i\) die Zahlen \(\frac12(n+1),\ \frac12(n+3),\ldots,n-1,\)
\end{tabular}\end{center}
\begin{center}\begin{tabular}{@{}l@{}}
\(f^0\) die Zahlen \(\frac{af}n\), deren ganzer Theil von der Form \(4k\)\\[3pt]
\(f'\) die Zahlen \(\frac{af}n\), deren ganzer Theil von der Form \(4k+1\)\\[3pt]
\([f''\) die Zahlen \(\frac{af}n\), deren ganzer Theil von der Form \(4k+2\)\\[3pt]
\(f'''\) die Zahlen \(\frac{af}n\), deren ganzer Theil von der Form \(4k+3]\)\\[3pt]
\(g^0\) die Zahlen \(\frac{ag}n\), deren ganzer Theil von der Form \(4k,\)
\end{tabular}\end{center}
\noindent [ebenso bezeichne \(g',g'',g'''\) die Zahlen \(\frac{ag}n\), deren ganzer Theil bezüglich von der Form \(4k+1,\ 4k+2,\ 4k+3\) ist und \(h^0,h',h'',h'''\) die Zahlen \(\frac{ah}n\), deren ganzer Theil bezüglich von der Form \(4k,4k+1,4k+2,4k+3\) ist, ferner \(i^0,i',i'',i'''\) die Zahlen \(\frac{ai}n\), deren ganzer Theil bezüglich von der Form \(4k,4k+1,4k+2,4k+3\) ist. Sei \(F^0,F',F'',F'''\) bezüglich die Anzahl der \(f^0,f',f'',f'''\), dieselbe
\newpage\SourceLeaf{79}\ArithmeticHead{72}{L}{BACHMANN, ÜBER GAUSS’ ZAHLENTHEORETISCHE ARBEITEN. ANHANG.}
\noindent Bedeutung sollen \(G^0,G',G'',G'''\) in Beziehung auf \(g\), \(H^0,H',H'',H'''\) in Beziehung auf \(h\), \(I^0,I',I'',I'''\) in Beziehung auf \(i\) haben. Dann ist:]
\[
\begin{aligned}
F^0+G^0&=H^0+I^0\\
F'+G'&=H'+I'\\
F''+G''&=H''+I''\\
F'''+G'''&=H'''+I'''.
\end{aligned}
\]
\begin{center}\begin{tabular}{@{}l@{}}
Jedes \(2h^0\) ist entweder ein \(g^0\) oder \(g'\)\\
Jedes \(2h'\) ist entweder ein \(g''\) oder \(g'''\)\\
Jedes \(2h''\) ist entweder ein \(g^0\) oder \(g'\)\\
Jedes \(2h'''\) ist entweder ein \(g''\) oder \(g''',\)
\end{tabular}\end{center}
\begin{center}\begin{tabular}{@{}rl@{}}
also & [die] \(2h^0\) vereinigt mit den \(2h''\) [giebt] \(g^0\) oder \(g'\)\\
& [die] \(2h'\) vereinigt mit den \(2h'''\) [giebt] \(g''\) oder \(g'''.\)
\end{tabular}\end{center}
\noindent\hspace*{\parindent}[Hieraus folgt]
\[
\begin{aligned}
H^0+H''&=G^0+G'\\
H'+H'''&=G''+G'''
\end{aligned}
\]
[und da \(F'''+G'''=H'''+I'''\)]
\[
\begin{aligned}
G'''&=H'''+I'''-F'''\\
&=H'+H'''-G''
\end{aligned}
\]
also
\[
I'''-F'''=H'-G''
\]
oder
\[
I'''-H'=F'''-G''.
\]
[Nun liegt] \(ai^0\) zwischen \(4kn\) und \(4kn+n\) [, also] \(a(n-i^0)\) zwischen \(an-(4kn+n)\) und \(an-4kn\).

[Ist] \(a=4\lambda+1\) [, so liegt] \(a(n-i^0)\) zwischen \(4(\lambda-k)n\) und \(4(\lambda-k)n+n\) [, d. h. es giebt so viel Zahlen \(i^0\) als \(h^0\). Ebenso findet man, dass es soviel Zahlen \(i'\) als \(h'''\), soviel Zahlen \(i''\) als \(h''\) und soviel Zahlen \(i'''\) als \(h'\) giebt. Also:]
\[
\begin{aligned}
I^0&=H^0\\
I'&=H'''\\
I''&=H''\\
I'''&=H'.
\end{aligned}
\]
\newpage\SourceLeaf{80}\ArithmeticHead{73}{R}{M. A. STERN, ÜBER EINEN SATZ VON GAUSS.}\HistoricalFoot{X\textsubscript{2}.}{10}
\noindent [Aus der letzten Gleichung folgt]
\[
F'''=G''.
\]
[Aus denselben Betrachtungen findet man aber auch]
\[
\begin{aligned}
F^0&=G^0\\
F'&=G'''\\
F''&=G''\\
F'''&=G',
\end{aligned}
\]
[also]
\[
F''=F''',
\]
[was zu beweisen war, und]
\[
G'=G''[\hbox{\footnotemark}].
\]
\begingroup\renewcommand{\thefootnote}{[*)}
\footnotetext{In dem Manuskripte steht] \(F'=F''',\ G'=G'''\).}
\endgroup

Ich bemerke noch schliesslich, dass ich diesen Satz auf ähnliche Weise in dem \textsc{Crelle}schen Journal für d. Mathematik Bd. 59 [, 1861, S. 146] bewiesen habe.

\begin{center}\small BEMERKUNGEN ZU DER VORSTEHENDEN NOTE.\end{center}
\begingroup\small
\noindent\hspace*{\parindent}I. Der von \textsc{Stern} mitgeteilte \textsc{Gauss}sche Satz steht in nächster Beziehung zu dem Werke X1, S 51 aus \textsc{Gauss}’ Nachlass abgedruckten Satze, der, soweit er sich auf ein ungerades \(n\) bezieht, leicht aus ihm ableitbar ist. Mit Beibehaltung der in der vorstehenden Note angewandten Bezeichnung und indem
\[
af=qn+r,\ 0\EqualAboveLess r<n;\quad q'=\left[\frac q4\right]
\]
gesetzt wird, entsprechen den Werten
\[
f=f^0,f',f'',f'''
\]
die Werte
\[
\begin{aligned}
\left[\frac{af}{4n}+\frac12\right]&=q',q',q'+1,q'+1\\
\left[\frac{af}{4n}\right]&=q',q',q',q'\\
\left[\frac{af}{4n}+\frac12\right]-\left[\frac{af}{4n}\right]&=0,0,1,1.
\end{aligned}
\]
Daher ist
\[
\Sigma_f\left(\left[\frac{af}{4n}+\frac12\right]-\left[\frac{af}{4n}\right]\right)=F''+F'''.
\]
Ebenso findet sich
\[
\Sigma_h\left(\left[\frac{ah}{4n}+\frac12\right]-\left[\frac{ah}{4n}+\frac14\right]\right)=H''.
\]
\newpage\SourceLeaf{81}\ArithmeticHead{74}{L}{BACHMANN, ÜBER GAUSS’ ZAHLENTHEORETISCHE ARBEITEN. ANHANG.}
\noindent Nach der vorstehenden Note ist aber
\[
F''+G''=H''+I'',
\]
andererseits
\[
G''=F''',\ H''=I'',
\]
also
\[
F''+F'''=2H'',
\]
wie in dem Satze Werke X1, S. 51 ausgesagt wird. Da \(F''=F'''\) ist, so folgt noch
\[
F''=F'''=H''.
\]
\begin{flushright}\textsc{Bachmann.}\end{flushright}

II. Angaben über Leben und Schriften von \textsc{Michel Reiss} (geb. 1805 zu Frankfurt a. M., gest. ebendaselbst 1869) findet man in der \textit{Allgemeinen Deutschen Biographie} XXVIII (1889), S. 143 und in \textsc{Poggendorffs} \textit{Biograph.-Literar. Handwörterbuch} III (1898), S. 1104. Die in der Einleitung zu der vorstehenden Note erwähnten Hefte hat der Sohn von M. A. \textsc{Stern}, Professor \textsc{Alfred Stern} (Zürich) nach dem Tode seines Vaters (1894) der Königl. Gesellschaft der Wissenschaften zu Göttingen geschenkt.
\begin{flushright}\textsc{Schlesinger.}\end{flushright}
\endgroup
\bigskip\begin{center}\rule{.24\textwidth}{.3pt}\end{center}
